\begin{table}[!htbp]\centering\tabsize
\caption{\Bi{Acquisition and selection counts by record type.}{按记录类型统计的获取与选择数量。}}\label{tab:acquisition}
\begin{tabularx}{\linewidth}{@{}p{0.24\linewidth}p{0.21\linewidth}Y@{}}\toprule
\Bi{Set}{集合} & \Bi{Count}{数量} & \Bi{Interpretation}{解释}\\\midrule
\Bi{Database hits / distinct records}{数据库命中／去重后记录} & 216/128; 75/47; 13/8; 10/7 & \Bi{Scopus; WoS; IEEE; ACM.}{依次为 Scopus、WoS、IEEE、ACM。}\\
\Bi{arXiv hits / distinct records}{arXiv 命中／去重后记录} & 544 / 346 & \Bi{API and keyword searches.}{API 与关键词检索。}\\
\Bi{Selection inventory}{选材清单} & 155 & \Bi{Candidate families used for framework source selection.}{用于框架来源选择的候选家族。}\\
\Bi{Selected / not selected}{采用／未采用} & 34 / 121 & \Bi{121 = 113 not used for claims + 8 contextual candidates.}{121 = 113 项未用于主张提取 + 8 项背景候选。}\\
\Bi{Coded sources outside inventory}{清单外编码来源} & 21 & \Bi{10 arXiv; 1 Scopus; 5 survey-collection; 5 targeted retrieval.}{10 项 arXiv、1 项 Scopus、5 项综述资料、5 项定向获取。}\\
\Bi{Coded subset}{编码子集} & 34 + 21 = 55 & \Bi{33 primary mechanisms and 22 auxiliary sources; 75 claims.}{33 个主要机制来源及 22 个辅助来源；75 条主张。}\\
\bottomrule\end{tabularx}
\end{table}
